% GENERATED FILE. Edit canonical lessons, then run npm run generate:books.
% canonical-source-hash: fnv1a64:39d60f7a68e5d915
% canonical-lessons: MR-C432-bhok, MR-C432-chidra, MR-C432-dag, MR-C432-viskatne, MR-C432-pajharne

\chapter{Holes and Stains}
\label{ch:mr-holes-and-stains}
\hlchaptermodality{432}

\begin{chapteropening}
\textbf{By the end of this chapter:} \emph{I can talk about holes and stains, and say things got messed up or oozed.}

Complete the last lesson of chapter 432: I can talk about holes and stains, and say things got messed up or oozed.
\end{chapteropening}

\section[\texorpdfstring{\mr{भोक} (bhok)}{भोक (bhok)}]{\mr{भोक} (bhok) — a hole}

\label{lesson:MR-C432-bhok}



\begin{quote}
Before the new one: say the Marathi for to burst, then the Marathi for to rip.
\end{quote}

\subsection*{You'll want to know: \mr{भोक}}
\textbf{\mr{भोक}} — \emph{bhok} — \textquotedblleft{}a hole\textquotedblright{}.

\begin{quote}
\mr{माझ्या} \mr{खिशाला} \mr{भोक} \mr{पडलं} \mr{आहे}.
\end{quote}

\begin{quote}
\emph{mājhyā khiśālā bhok paḍlaṁ āhe.}
\end{quote}

\begin{quote}
There is a hole in my pocket.
\end{quote}

Neuter.

\mr{भोक} \mr{पडणे}, a hole appears; the thing with the hole takes -\mr{ला}.

\begin{cousinweb}
\textbf{\mr{भोक}} (\emph{bhok}) is inherited: CDIAL lists Marathi \emph{bhok}, \textquotedblleft{}hole\textquotedblright{}, under \emph{*bhōkka}, \textquotedblleft{}hollow\textquotedblright{} (CDIAL 9624).
\end{cousinweb}

\begin{grammarlens}[title={What you've built}]
One more word for getting around.
\end{grammarlens}

\subsection*{Your turn}
\begin{itemize}
\raggedright
  \item \emph{Say it:} \emph{bhok}
  \item \emph{Say it:} \emph{bhok}, once more
  \item \emph{Say it:} \emph{mājhyā khiśālā bhok paḍlaṁ āhe.} — There is a hole in my pocket.
  \item \emph{Recall:} say the Marathi for to burst, then the Marathi for to rip, then say \emph{bhok} again
\end{itemize}

\begin{tcolorbox}[breakable,colback=teal!4,colframe=teal!35!black,title={Before you move on}]
What does \mr{भोक} mean? (\textquotedblleft{}A hole\textquotedblright{}.) Say it once more.
\end{tcolorbox}
\section[\texorpdfstring{\mr{छिद्र} (chidra)}{छिद्र (chidra)}]{\mr{छिद्र} (chidra) — a small hole; a pore}

\label{lesson:MR-C432-chidra}



\begin{quote}
Before the new one: say the Marathi for to rip, then the Marathi for a hole.
\end{quote}

\subsection*{You'll want to know: \mr{छिद्र}}
\textbf{\mr{छिद्र}} — \emph{chidra} — \textquotedblleft{}a small hole; a pore\textquotedblright{}.

\begin{quote}
\mr{बादलीच्या} \mr{तळाला} \mr{एक} \mr{छिद्र} \mr{आहे}.
\end{quote}

\begin{quote}
\emph{bādlīcyā taḷālā ek chidra āhe.}
\end{quote}

\begin{quote}
There is a small hole in the bottom of the bucket.
\end{quote}

Neuter.

More formal and smaller than \mr{भोक}.

\begin{cousinweb}
\textbf{\mr{छिद्र}} (\emph{chidra}) is a tatsama, taken from Sanskrit \emph{chidra}, \textquotedblleft{}torn open; a hole\textquotedblright{} (Monier-Williams).
\end{cousinweb}

\begin{grammarlens}[title={What you've built}]
One more word for getting around.
\end{grammarlens}

\subsection*{Your turn}
\begin{itemize}
\raggedright
  \item \emph{Say it:} \emph{chidra}
  \item \emph{Say it:} \emph{chidra}, once more
  \item \emph{Say it:} \emph{bādlīcyā taḷālā ek chidra āhe.} — There is a small hole in the bottom of the bucket.
  \item \emph{Recall:} say the Marathi for to rip, then the Marathi for a hole, then say \emph{chidra} again
\end{itemize}

\begin{tcolorbox}[breakable,colback=teal!4,colframe=teal!35!black,title={Before you move on}]
What does \mr{छिद्र} mean? (\textquotedblleft{}A small hole; a pore\textquotedblright{}.) Say it once more.
\end{tcolorbox}
\section[\texorpdfstring{\mr{डाग} (ḍāg)}{डाग (ḍāg)}]{\mr{डाग} (ḍāg) — a stain; a spot}

\label{lesson:MR-C432-dag}



\begin{quote}
Before the new one: say the Marathi for a hole, then the Marathi for a small hole.
\end{quote}

\subsection*{You'll want to know: \mr{डाग}}
\textbf{\mr{डाग}} — \emph{ḍāg} — \textquotedblleft{}a stain; a spot\textquotedblright{}.

\begin{quote}
\mr{शर्टवर} \mr{चहाचा} \mr{डाग} \mr{पडला}.
\end{quote}

\begin{quote}
\emph{śarṭvar cahācā ḍāg paḍlā.}
\end{quote}

\begin{quote}
A tea stain got on the shirt.
\end{quote}

Masculine.

\mr{डाग} \mr{पडणे}, a stain appears; \mr{डाग} \mr{काढणे}, to remove one.

\begin{cousinweb}
\textbf{\mr{डाग}} (\emph{ḍāg}) is a loan, usually traced to Persian \emph{dāġ}, \textquotedblleft{}a brand, a mark\textquotedblright{}.
\end{cousinweb}

\begin{grammarlens}[title={What you've built}]
One more word for getting around.
\end{grammarlens}

\subsection*{Your turn}
\begin{itemize}
\raggedright
  \item \emph{Say it:} \emph{ḍāg}
  \item \emph{Say it:} \emph{ḍāg}, once more
  \item \emph{Say it:} \emph{śarṭvar cahācā ḍāg paḍlā.} — A tea stain got on the shirt.
  \item \emph{Recall:} say the Marathi for a hole, then the Marathi for a small hole, then say \emph{ḍāg} again
\end{itemize}

\begin{tcolorbox}[breakable,colback=teal!4,colframe=teal!35!black,title={Before you move on}]
What does \mr{डाग} mean? (\textquotedblleft{}A stain; a spot\textquotedblright{}.) Say it once more.
\end{tcolorbox}
\section[\texorpdfstring{\mr{विस्कटणे} (viskaṭṇe)}{विस्कटणे (viskaṭṇe)}]{\mr{विस्कटणे} (viskaṭṇe) — to get messed up}

\label{lesson:MR-C432-viskatne}



\begin{quote}
Before the new one: say the Marathi for a small hole, then the Marathi for a stain.
\end{quote}

\subsection*{You'll want to know: \mr{विस्कटणे}}
\textbf{\mr{विस्कटणे}} — \emph{viskaṭṇe} — \textquotedblleft{}to get messed up\textquotedblright{}.

\begin{quote}
\mr{धावताना} \mr{तिचे} \mr{केस} \mr{विस्कटले}.
\end{quote}

\begin{quote}
\emph{dhāvtānā tice kes viskaṭle.}
\end{quote}

\begin{quote}
Her hair got messed up as she ran.
\end{quote}

Past (he …): \mr{तो} \mr{विस्कटला}, \emph{to viskaṭlā}.

Hair, papers and plans can all \mr{विस्कटणे}.

\begin{cousinweb}
\textbf{\mr{विस्कटणे}} (\emph{viskaṭṇe}) is inherited: CDIAL lists Marathi \emph{viskaṭṇẽ}, \textquotedblleft{}to disorder, derange\textquotedblright{}, under \emph{*viṣvaskati}, \textquotedblleft{}goes astray\textquotedblright{} (CDIAL 11995).
\end{cousinweb}

\begin{grammarlens}[title={What you've built}]
One more word for getting around.
\end{grammarlens}

\subsection*{Your turn}
\begin{itemize}
\raggedright
  \item \emph{Say it:} \emph{viskaṭṇe}
  \item \emph{Say it:} \emph{viskaṭṇe}, once more
  \item \emph{Say it:} \emph{dhāvtānā tice kes viskaṭle.} — Her hair got messed up as she ran.
  \item \emph{Recall:} say the Marathi for a small hole, then the Marathi for a stain, then say \emph{viskaṭṇe} again
\end{itemize}

\begin{tcolorbox}[breakable,colback=teal!4,colframe=teal!35!black,title={Before you move on}]
What does \mr{विस्कटणे} mean? (\textquotedblleft{}To get messed up\textquotedblright{}.) Say it once more.
\end{tcolorbox}
\section[\texorpdfstring{\mr{पाझरणे} (pājharṇe)}{पाझरणे (pājharṇe)}]{\mr{पाझरणे} (pājharṇe) — to ooze}

\label{lesson:MR-C432-pajharne}



\begin{quote}
Before the new one: say the Marathi for a stain, then the Marathi for to get messed up.
\end{quote}

\subsection*{You'll want to know: \mr{पाझरणे}}
\textbf{\mr{पाझरणे}} — \emph{pājharṇe} — \textquotedblleft{}to ooze\textquotedblright{}.

\begin{quote}
\mr{भिंतीतून} \mr{पाणी} \mr{पाझरत} \mr{होतं}.
\end{quote}

\begin{quote}
\emph{bhintītūn pāṇī pājhrat hotaṁ.}
\end{quote}

\begin{quote}
Water was oozing through the wall.
\end{quote}

Past (he …): \mr{तो} \mr{पाझरला}, \emph{to pājharlā}.

The source takes -\mr{तून}: \mr{भिंतीतून}, through the wall.

\begin{cousinweb}
\textbf{\mr{पाझरणे}} (\emph{pājharṇe}) is inherited: CDIAL lists Marathi \emph{pājharṇẽ}, \textquotedblleft{}to drip down\textquotedblright{}, under \emph{*prajjharati}, \textquotedblleft{}flows\textquotedblright{} (CDIAL 8512).
\end{cousinweb}

\begin{grammarlens}[title={What you've built}]
One more word for getting around.
\end{grammarlens}

\subsection*{Your turn}
\begin{itemize}
\raggedright
  \item \emph{Say it:} \emph{pājharṇe}
  \item \emph{Say it:} \emph{pājharṇe}, once more
  \item \emph{Say it:} \emph{bhintītūn pāṇī pājhrat hotaṁ.} — Water was oozing through the wall.
  \item \emph{Recall:} say the Marathi for a stain, then the Marathi for to get messed up, then say \emph{pājharṇe} again
\end{itemize}

\begin{tcolorbox}[breakable,colback=teal!4,colframe=teal!35!black,title={Before you move on}]
What does \mr{पाझरणे} mean? (\textquotedblleft{}To ooze\textquotedblright{}.) Say it once more.
\end{tcolorbox}
